%% PPgEA_UFRN - Tutorial
% --------------------------------------------------------------------
%% Modelo de dissertação adaptado de
%% abtex2-modelo-trabalho-academico.tex, v-1.9.7 laurocesar
%% Copyright 2012-2018 by abnTeX2 group at http://www.abntex.net.br/ 
%%

\documentclass[
	% -- opções da classe memoir --
	12pt,				% tamanho da fonte
	openany,			% evita páginas em branco para iniciar capítulos em página ímpar
	%openright,			% capítulos começam em pág ímpar (insere página vazia caso preciso)
	twoside,			% para impressão em recto e verso. Oposto a oneside
	%oneside,
	a4paper,			% tamanho do papel. 
	% -- opções da classe abntex2 --
	%chapter=TITLE,		% títulos de capítulos convertidos em letras maiúsculas
	%section=TITLE,		% títulos de seções convertidos em letras maiúsculas
	%subsection=TITLE,	% títulos de subseções convertidos em letras maiúsculas
	%subsubsection=TITLE,% títulos de subsubseções convertidos em letras maiúsculas
	% -- opções do pacote babel --
	english,			% idioma adicional para hifenização
	french,				% idioma adicional para hifenização
	spanish,			% idioma adicional para hifenização
	brazilian			% o último idioma é o principal do documento
	]{abntex2}

% --------------------------------------------------------------------
% Pacotes básicos 
% --------------------------------------------------------------------
\usepackage[T1]{fontenc}		% Selecao de codigos de fonte.
\usepackage[utf8]{inputenc}		% Codificacao do documento (conversão automática dos acentos)
\usepackage{indentfirst}		% Indenta o primeiro parágrafo de cada seção.
\usepackage{color}				% Controle das cores
\usepackage{graphicx}			% Inclusão de gráficos
\usepackage{microtype} 			% para melhorias de justificação

% --------------------------------------------------------------------
% Pacotes adicionais, usados apenas no âmbito do Modelo Canônico do abnteX2
% --------------------------------------------------------------------
% \usepackage{lipsum}				% para geração de dummy text

\usepackage{PPgEA}
% ---------------------------------------------------------------------
% Informações de dados para CAPA e FOLHA DE ROSTO
% ---------------------------------------------------------------------
\titulo{Modelo de Dissertação PPgEA com \abnTeX}
\autor{Autor}
\local{Natal}
\data{\the\year}
\orientador{Orientador}
\coorientador{Coorientador}
\instituicao{%
  Universidade Federal do Rio Grande do Norte\\
  %\par
  Escola de Ciências e Tecnologia\\
  %\par
  Programa de Pós-Graduação em Engenharia Aeroespacial}
  

\tipotrabalho{Dissertação}

% O preambulo deve conter o tipo do trabalho, o objetivo, 
% o nome da instituição e a área de concentração 
\preambulo{Dissertação apresentada ao Programa de Pós-Graduação em Engenharia Aeroespacial da Universidade Federal do Rio Grande do Norte, como parte dos requisitos para a obtenção do título de \textbf{Mestre em Engenharia Aeroespacial}, área de concentração: \textbf{Materiais e Tecnologias Aeroespaciais}}

% ---------------------------------------------------------------------


 % CAPA e FOLHA DE ROSTO

%--------------------------------------------------------------
% compila o indice
%--------------------------------------------------------------
\makeindex

%--------------------------------------------------------------
% Início do documento
%--------------------------------------------------------------
\begin{document}

% Seleciona o idioma do documento (conforme pacotes do babel)
%\selectlanguage{english}
\selectlanguage{brazilian}
\frenchspacing      % Retira espaço extra obsoleto entre as frases.

% ----------------------------------------------------------
% ELEMENTOS PRÉ-TEXTUAIS
% ----------------------------------------------------------
% \pretextual

% Capa
\imprimircapa

% Folha de rosto - (o * indica que haverá a ficha bibliográfica)
\imprimirfolhaderosto*

% ---------------------------------------------------------------------
% Ficha bibliografica
% ---------------------------------------------------------------------
\begin{fichacatalografica}
	\sffamily
	\vspace*{\fill}					% Posição vertical
	\begin{center}					% Minipage Centralizado
		\fbox{\begin{minipage}[c][8cm]{13.3cm}		% Largura
		\small
		\raggedright
	\imprimirautor
	%Sobrenome, Nome do autor
	
	\hspace{0.5cm} \imprimirtitulo  / \imprimirautor. --
	\imprimirlocal, \imprimirdata-
	
	\hspace{0.5cm} \thelastpage p.\\
	
	\hspace{0.5cm} \imprimirorientadorRotulo~\imprimirorientador\\ 

	\hspace{0.5cm} \imprimircoorientadorRotulo~\imprimircoorientador\\
	
		\hspace{0.5cm}
		\parbox[t]{\dimexpr\linewidth-0.5cm\relax}{\imprimirtipotrabalho~--~\imprimirinstituicao,
		\imprimirdata.}\\
		
		\hspace{0.5cm}\parbox[t]{\dimexpr\linewidth-0.5cm\relax}{%
			1. Palavra-chave1.
			2. Palavra-chave2.
			3. Palavra-chave3.
			I. Orientador.
			II. Universidade xxx.
			III. Faculdade de xxx.
			IV. Título.}			
		\end{minipage}}
	\end{center}
\end{fichacatalografica}
 
% A Biblioteca Central da UFRN fornecerá um PDF com a ficha catalográfica definitiva após a defesa do trabalho. Quando estiver com o documento, salve-o como PDF no diretório do seu projeto e substitua todo o conteúdo de implementação deste arquivo pelo comando abaixo:

% \begin{fichacatalografica}
%     \includepdf{fig_ficha_catalografica.pdf}
% \end{fichacatalografica}

% ---------------------------------------------------------------------
% Inserir errata
% ---------------------------------------------------------------------

%\input{errata}

% ---------------------------------------------------------------------
% Inserir folha de aprovação
% ---------------------------------------------------------------------
% Elemento obrigatório da NBR 14724:2024 (seção correspondente à folha de aprovação). 
\begin{folhadeaprovacao}

  \begin{center}
    {\ABNTEXchapterfont\large\imprimirautor}

    \vspace*{\fill}\vspace*{\fill}
    \begin{center}
      \ABNTEXchapterfont\bfseries\Large\imprimirtitulo
    \end{center}
    \vspace*{\fill}
    
    \hspace{.45\textwidth}
    \begin{minipage}{.5\textwidth}
        \imprimirpreambulo
    \end{minipage}%
    \vspace*{\fill}
   \end{center}
        
   Trabalho aprovado.  \imprimirlocal, \dataDefesa:

   \assinatura{\textbf{\imprimirorientador} \\ Orientador}
   \assinatura{\textbf{\bancaConvidadoUm} \\ Convidado 1}
   \assinatura{\textbf{\bancaConvidadoDois} \\ Convidado 2}
   %\assinatura{\textbf{\bancaConvidadoTres} \\ Convidado 3}
   %\assinatura{\textbf{\bancaConvidadoQuatro} \\ Convidado 4}
      
   \begin{center}
    \vspace*{0.5cm}
    {\large\imprimirlocal}
    \par
    {\large\imprimirdata}
    \vspace*{1cm}
  \end{center}
  
\end{folhadeaprovacao}
% Você pode utilizar este modelo até a aprovação
% do trabalho. Após isso, substitua todo o conteúdo deste arquivo por uma
% imagem da página assinada pela banca com o comando abaixo:
%
% \begin{folhadeaprovacao}
% \includepdf{folhadeaprovacao_final.pdf}
% \end{folhadeaprovacao}
%

% ---------------------------------------------------------------------
% Dedicatória
% ---------------------------------------------------------------------
\input{elementos_pre_textuais/dedicatoria}

% ---------------------------------------------------------------------
% Agradecimentos
% ---------------------------------------------------------------------
\begin{agradecimentos}
Espaço destinado aos agradecimentos do autor.

\end{agradecimentos}

% ---------------------------------------------------------------------
% Epígrafe
% ---------------------------------------------------------------------
\input{elementos_pre_textuais/epigrafe}

% ---------------------------------------------------------------------
% RESUMOS
% ---------------------------------------------------------------------
\input{elementos_pre_textuais/resumo_abstract}

% ---------------------------------------------------------------------
% Lista de ilustrações
% ---------------------------------------------------------------------
\pdfbookmark[0]{\listfigurename}{lof}
\listoffigures*
\cleardoublepage

% ---------------------------------------------------------------------
% Lista de quadros
% ---------------------------------------------------------------------
\pdfbookmark[0]{\listofquadrosname}{loq}
\listofquadros*
\cleardoublepage

% ---------------------------------------------------------------------
% Lista de tabelas
% ---------------------------------------------------------------------
\pdfbookmark[0]{\listtablename}{lot}
\listoftables*
\cleardoublepage

% ---------------------------------------------------------------------
% Lista de abreviaturas e siglas
% ---------------------------------------------------------------------
\input{elementos_pre_textuais/lista_abrev}

% ---------------------------------------------------------------------
% Lista de símbolos
% ---------------------------------------------------------------------
\input{elementos_pre_textuais/lista_simb}

% ---------------------------------------------------------------------
% Sumário
% ---------------------------------------------------------------------
\pdfbookmark[0]{\contentsname}{toc}
\tableofcontents*
\cleardoublepage

% ----------------------------------------------------------
% ELEMENTOS TEXTUAIS - CAPÍTULOS
% ----------------------------------------------------------
\textual

\include{elementos_textuais/intro_cap}
\include{elementos_textuais/abntex2-modelo-include-comandos}
\include{elementos_textuais/3_cap}
\include{elementos_textuais/4_cap}
\include{elementos_textuais/5_cap}

% ----------------------------------------------------------
% ELEMENTOS PÓS-TEXTUAIS
% ----------------------------------------------------------
\postextual
% ----------------------------------------------------------

% ----------------------------------------------------------
% Referências bibliográficas
% ----------------------------------------------------------
\bibliography{referencias}

% ----------------------------------------------------------
% Glossário
% ----------------------------------------------------------
%
% Consulte o manual da classe abntex2 para orientações sobre o glossário.
%
%\glossary

% ----------------------------------------------------------
% Apêndices
% ----------------------------------------------------------

\begin{apendicesenv}
    % Imprime uma página indicando o início dos apêndices
    \partapendices
    \include{elementos_pos_textuais/apend_A}
    % ----------------------------------------------------------
\chapter{Nullam elementum urna vel imperdiet eu metus}
% ----------------------------------------------------------
\lipsum[55-57]
\end{apendicesenv}

% ----------------------------------------------------------
% Anexos
% ----------------------------------------------------------
\begin{anexosenv}
  % Imprime uma página indicando o início dos anexos
    \partanexos
    \include{elementos_pos_textuais/anex_A}
    % ---
\chapter{Fusce facilisis lacinia dui}
% ---

\lipsum[30]

    \include{elementos_pos_textuais/anex_C}

\end{anexosenv}

%---------------------------------------------------------------------
% INDICE REMISSIVO
%---------------------------------------------------------------------
\phantompart
\printindex


\end{document}
